% Pista pro-24i-q4 · Probabilitat
% Procedència: PAU Catalunya 2024 · Convocatòria de juny (incidències) · Sèrie 5 · Qüestió 4
\documentclass[11pt,a4paper]{article}
\usepackage{pista}
\pistainfo{Catalunya 2024}{Convocatòria de juny (incidències)}{Sèrie 5}

\begin{document}

\section*{\textcolor{blauPista}{Qüestió 4}}

\noindent La Rut fa servir el mètode següent per a fer els problemes de matemàtiques: tira un dau equilibrat i, si el resultat és com a màxim 4, pensa i resol el problema ella mateixa; si el resultat és 5 o 6, busca la solució del problema per Internet i la copia. Quan és ella qui ha pensat la solució, la resposta és correcta en el $75\,\%$ dels casos; quan copia la solució d'Internet, la resposta és correcta només en el $40\,\%$ dels casos.

\subsection*{\textit{a)} \normalfont Quina és la probabilitat que la solució d'un problema respost seguint aquest mètode sigui correcta? \hspace*{\fill}\textnormal{[0,75~punts]}}

\pista{Fes un diagrama d'arbre. Hi ha dos camins que porten a una resposta correcta:
\begin{itemize}
\item la Rut pensa la solució \emph{i} és correcta;
\item la Rut la copia d'Internet \emph{i} és correcta.
\end{itemize}
Aplica el teorema de la probabilitat total: multiplica al llarg de cada camí i suma els dos camins.

\medskip
\textbf{Atenció:} les probabilitats de la primera branca surten del dau: «com a màxim 4» vol dir 1, 2, 3 o 4.}

\subsection*{\textit{b)} \normalfont Quina és la probabilitat que un problema l'hagi resolt la Rut si sabem que la solució és correcta? \hspace*{\fill}\textnormal{[0,75~punts]}}

\pista{És un cas clar del teorema de Bayes: et demanen una probabilitat condicionada en sentit \emph{invers} al de les dades.

\medskip
Anomena $R$ = «l'ha resolt la Rut» i $C$ = «la solució és correcta».
\begin{itemize}
\item Coneixes $P(C \mid R)$: «és correcta \emph{sabent que} l'ha resolt ella».
\item Et demanen $P(R \mid C)$: «l'ha resolt ella \emph{sabent que} és correcta».
\end{itemize}
El denominador, $P(C)$, és el resultat de l'apartat \textit{a)}.}

\subsection*{\textit{c)} \normalfont Demà la Rut ha d'entregar 5 problemes de matemàtiques. Quina és la probabilitat que n'hi hagi almenys 4 de correctes? \hspace*{\fill}\textnormal{[1~punt]}}

\pista{Sigui $X$ el nombre de problemes correctes dels 5. Aleshores $X \sim B(5, p)$, on $p = P(C)$ és el resultat de l'apartat \textit{a)}.

\medskip
«Almenys 4 correctes» vol dir «exactament 4 o exactament 5»:
\[
P(X \ge 4) = P(X = 4) + P(X = 5).
\]}

\end{document}
